\section{Elementary Gaussian constants}\label{ms:normalization}\label{ms:gaussian-constant-proof}
\begin{lemma}[Elementary Gaussian constants]\label{ms:gaussian-constants}
For the standard normal density,
\[
 \|\phi''\|_\infty=\phi(0)<2/5,\qquad
 \|\phi'''\|_\infty\le2/\pi,\qquad 1/(4\pi)<2/25.
\]
Moreover $819/2000<\CE<20487/50000<41/100$ and
$\phi(0)>0.39886>0.3988=997/2500$.
\end{lemma}
We give the exact comparisons used in Lemma~\ref{ms:gaussian-constants}.

\begin{proof}
Since $\phi''(x)=(x^2-1)\phi(x)$, the norm on $x^2\le1$ is at most
$\phi(0)$. On $x^2\ge1$, the maximum of
$(x^2-1)e^{-x^2/2}$ is $2e^{-3/2}<1$, attained when $x^2=3$.
Thus the norm is exactly $\phi(0)$.

For completeness, the Gaussian transform can be obtained without an
inversion theorem for the distribution being studied. Differentiating
$J(t)=\int_\R e^{-x^2/2}e^{itx}\,dx$ under the integral and integrating by
parts give $J'(t)=-tJ(t)$. The Gaussian integral, evaluated by squaring it
and using polar coordinates, gives $J(0)=\sqrt{2\pi}$. Hence
$J(t)=\sqrt{2\pi}e^{-t^2/2}$; applying the same formula with the variables
interchanged gives
\[
 \phi(x)=\frac1{2\pi}\int_\R e^{-t^2/2}e^{-itx}\,dt.
\]
Three differentiations are justified by the integrability of
$|t|^3e^{-t^2/2}$. Consequently
\[
 \|\phi'''\|_\infty
 \le\frac1{2\pi}\int_\R|t|^3e^{-t^2/2}\,dt=\frac2\pi.
\]
The last integral equals four after the substitution $w=t^2/2$ on the
positive half-line.

Here are elementary rational comparisons for all numerical constants used
in this section. Put $a=\arctan(1/5)$ and $b=\arctan(1/239)$.
Tangent addition gives $\tan(2a)=5/12$, $\tan(4a)=120/119$ and
$\tan(4a-b)=1$. Also
\[
 0<4(1/5-1/(3\cdot5^3))-1/239<4a-b<4/5<1<\pi/2.
\]
The arctangent inequalities follow by integrating
$1-t^2<1/(1+t^2)<1$, while $\pi/2=2\arctan(1)>1$ follows by integrating
$1/(1+t^2)>1/2$ on $[0,1)$. It follows that $\pi=16a-4b$ and
\[
 \pi>16(1/5-1/(3\cdot5^3))-4/239
      =281476/89625>25/8.
\]
This proves $\phi(0)<2/5$ and $1/(4\pi)<2/25$.

To bound $\CE$ more closely, define
\[
 s_k(x)=\sum_{j=0}^{k-1}\frac{(-1)^jx^{2j+1}}{2j+1},\quad
 \pi_-=16s_6(1/5)-4s_1(1/239),\quad
 \pi_+=16s_5(1/5)-4s_2(1/239).
\]
The remainder in the finite geometric expansion of $1/(1+t^2)$ has its
stated alternating sign, so integration gives $\pi_-<\pi<\pi_+$.
Let $r_-=3162277660/10^9$ and $r_+=3162277661/10^9$. Squaring shows
$r_-^2<10<r_+^2$. The following comparisons are exact rational ones:
\[
 \left(\frac{819}{2000}\right)^2
 <\frac{19+6r_-}{72\pi_+}
 <\CE^2=\frac{19+6\sqrt{10}}{72\pi}
 <\frac{19+6r_+}{72\pi_-}
 <\left(\frac{20487}{50000}\right)^2.
\]
For example, each outer comparison follows by multiplying by its positive
rational denominator; the displayed finite sums specify those denominators
completely. All quantities are positive, so taking square roots proves the
claimed bounds. In particular, $0.4095<\CE<0.40974<0.41$.
The same exact endpoints also satisfy
\[
 \pi_->\frac{333}{106}>\left(\frac{4431}{2500}\right)^2,
 \qquad \pi_+<\frac{22}{7}<2\left(\frac{251}{200}\right)^2.
\]
Thus $\sqrt\pi>4431/2500$ and $\sqrt{\pi/2}<251/200$,
as used in the Fourier bounds.
Since $(0.39886)^2<7/44$, we also obtain
$\phi(0)>0.39886>0.3988$.
\end{proof}

\begin{lemma}[The narrow entry enclosure]\label{ms:narrow-constant}
The target constant satisfies
\[
 c_-:=\frac{4097321837023}{10^{13}}<\CE
 <\frac{4097321837025}{10^{13}}=:c_+.
\]
\end{lemma}
\begin{proof}
Use the same finite sums $s_k$ and Machin identity proved above, now with
\[
 \pi_L=16s_{10}(1/5)-4s_5(1/239),\qquad
 \pi_U=16s_9(1/5)-4s_4(1/239).
\]
An even number of terms bounds each arctangent from below and an odd
number bounds it from above. The negative coefficient of the second
arctangent reverses its choice of endpoint, so
$0<\pi_L<\pi<\pi_U$. Put
\[
 r_L=\frac{3162277660168379}{10^{15}},\qquad
 r_U=\frac{3162277660168380}{10^{15}}.
\]
Integer squaring gives $r_L^2<10<r_U^2$. Direct rational comparisons give
\[
 c_-^2<\frac{19+6r_L}{72\pi_U}
 <\CE^2
 <\frac{19+6r_U}{72\pi_L}<c_+^2.
\]
The finite sums specify every rational in this display; each outer
inequality can be checked by cross-multiplying positive integers.
All terms are positive, so taking square roots gives the assertion.
This directed enclosure, rather than the coarser bounds in the preceding
lemma, supplies the precision used by the entry calculation.
\end{proof}
